\documentclass[11pt]{article}

\usepackage[margin=1in]{geometry}
\usepackage[T1]{fontenc}
\usepackage[utf8]{inputenc}
\usepackage{lmodern}
\usepackage{microtype}
\usepackage{xcolor}
\usepackage{hyperref}
\usepackage{booktabs}
\usepackage{longtable}
\usepackage{array}
\usepackage{enumitem}
\usepackage{verbatim}

\hypersetup{
  colorlinks=true,
  linkcolor=blue!50!black,
  urlcolor=blue!50!black
}

\setlength{\parindent}{0pt}
\setlength{\parskip}{0.6em}
\setlist[itemize]{leftmargin=1.5em}
\setlist[enumerate]{leftmargin=1.7em}
\sloppy

\newcolumntype{L}[1]{>{\raggedright\arraybackslash}p{#1}}
\newcommand{\fname}[1]{\path{#1}}

\title{\textbf{QASPER RAG Project Walkthrough}}
\author{CS505 NLP Final Project Team}
\date{\today}

\begin{document}
\maketitle

\section{Project Overview}
This project builds a \textbf{local, reproducible retrieval-augmented question answering system} for scientific papers using the \textbf{QASPER} dataset. The task is: given a paper and a question about that paper, retrieve the most relevant evidence passages and generate a grounded answer. In the citation-forcing setting, the system must also cite the retrieved evidence it claims to rely on.

The system solves three linked subproblems:
\begin{enumerate}
  \item \textbf{Preprocessing and chunking:} convert raw scientific papers into stable, retrievable passages.
  \item \textbf{Retrieval:} find the most relevant chunks for each question.
  \item \textbf{Generation and grounding:} generate an answer, optionally with citations, and evaluate whether the answer is correct and whether the citations are actually supported.
\end{enumerate}

At a high level, the project:
\begin{itemize}
  \item preprocesses QASPER into canonical shared artifacts,
  \item supports fixed, sliding-window, and section-aware chunking,
  \item evaluates BM25, dense retrieval, hybrid retrieval, and reranking,
  \item generates answers with local models such as Phi and Qwen,
  \item evaluates answer quality, semantic similarity, grounding, and citation quality,
  \item produces report-ready tables, analyses, and human-evaluation artifacts.
\end{itemize}

\section{Repository Structure}
\subsection{Top-Level Layout}
\begin{verbatim}
NLP_research_project/
  src/qasper_rag/              core library code
  scripts/                     entrypoints, batch jobs, backfills, analysis
  tests/                       unit and regression tests
  reports/                     proposal, midway report, status report, artifacts
  local_scc_results/           archived SCC outputs copied locally
  README.md                    repo usage guide
  setup_env.sh                 environment setup
  environment.yml              conda environment spec
\end{verbatim}

\subsection{What Each Folder Is For}
\begin{itemize}
  \item \texttt{src/qasper\_rag/}: the actual implementation of loading, chunking, retrieval, generation, and evaluation.
  \item \texttt{scripts/}: operational entrypoints used to build data, run experiments, submit SCC jobs, backfill metrics, analyze errors, and prepare human evaluation.
  \item \texttt{tests/}: correctness checks for the main components and regression tests for tricky bugs.
  \item \texttt{reports/}: LaTeX papers and supporting artifacts used for the proposal, midway report, and internal status tracking.
  \item \texttt{local\_scc\_results/}: a persistent local mirror of important SCC outputs so experiment JSONs are not lost.
\end{itemize}

\section{File-by-File Guide}
\subsection{Core Library Files}
\scriptsize
\begin{longtable}{L{0.15\textwidth} L{0.24\textwidth} L{0.27\textwidth} L{0.26\textwidth}}
\toprule
\textbf{File} & \textbf{Role} & \textbf{Key functions/classes} & \textbf{Connections} \\
\midrule
\endhead
\fname{config.py} &
Central path and layout resolver. &
\fname{ChunkingSpec}, \fname{ProjectPaths}, \fname{get_default_project_paths()}, \fname{ensure_project_layout()}. &
Used by almost every script to resolve SCC shared paths, per-user run paths, and local fallbacks. \\

\fname{schema.py} &
Normalized in-memory representation of raw QASPER. &
\fname{QasperPaper}, \fname{QasperQuestion}, \fname{QasperAnswer}, \fname{Section}, \fname{TextUnit}, \fname{Chunk}. &
Produced by \texttt{loader.py}; consumed by \texttt{chunking.py} and \texttt{processing.py}. \\

\fname{loader.py} &
Loads raw QASPER files or Hugging Face dataset records. &
\fname{load_qasper_json()}, \fname{load_qasper_huggingface()}, \fname{resolve_qasper_split_path()}, \fname{parse_qasper_paper()}. &
Feeds normalized paper/question/answer objects into \texttt{processing.py}. \\

\fname{chunking.py} &
Implements all chunking strategies. &
\fname{BaseChunker}, \fname{FixedChunker}, \fname{SlidingWindowChunker}, \fname{SectionAwareChunker}, token-window helpers. &
Called by \texttt{processing.py} to build retrievable chunk corpora. \\

\fname{processing.py} &
Builds canonical processed artifacts from raw QASPER. &
\fname{build_processed_split()}, \fname{process_qasper_split_file()}, record serializers and JSONL writers. &
Bridges raw input and the shared \fname{processed/} interface used by the rest of the pipeline. \\

\fname{artifacts.py} &
Loads processed artifacts back from disk. &
\fname{ProcessedQuestion}, \fname{ProcessedChunk}, \fname{load_processed_questions()}, \fname{load_processed_chunks()}, \fname{group_chunks_by_paper()}. &
Used by retrieval and generation runners. \\

\fname{retrieval.py} &
Implements the retrieval stack. &
\fname{BM25Retriever}, \fname{DenseRetriever}, \fname{HybridRetriever}, \fname{CrossEncoderReranker}, \fname{RerankRetriever}, \fname{NullRetriever}, \fname{RandomRetriever}, \fname{build_retriever()}. &
Used by both retrieval-only and retrieval+generation experiments. \\

\fname{retrieval_eval.py} &
Evaluates retrieval quality against QASPER evidence. &
\fname{QuestionRetrievalMetrics}, \fname{evaluate_ranked_chunks()}, \fname{aggregate_question_metrics()}, evidence matching helpers. &
Used by \fname{evaluate_qasper_retrieval.py}. \\

\fname{generation.py} &
Builds prompts, runs local generation, parses citations, normalizes outputs. &
\fname{PromptPackage}, \fname{GenerationPrediction}, \fname{HuggingFaceGenerator}, \fname{build_prompt()}, \fname{prepare_prompt_chunks()}, \fname{normalize_generation_answer()}, \fname{build_generation_prediction()}. &
Central generation logic used by the generation runner. \\

\fname{generation_eval.py} &
Evaluates answer quality and grounding. &
\fname{QuestionGenerationMetrics}, \fname{token_f1_score()}, \fname{collect_gold_answers()}, \fname{evaluate_generation_prediction()}, \fname{aggregate_generation_metrics()}. &
Used during generation experiments and in metric backfills. \\

\fname{nli_eval.py} &
Checks whether cited claims are supported by evidence. &
\fname{HuggingFaceEntailmentScorer}, \fname{CitedClaim}, \fname{extract_cited_claims()}, \fname{evaluate_cited_claims()}. &
Used for NLI citation backfill and citation verification. \\

\fname{citation_verifier.py} &
New verifier gate for citation-forcing outputs. &
\fname{verify_citation_prediction()}. &
Called inside the generation runner to keep only supported cited claims or abstain to \fname{unanswerable}. \\

\fname{__init__.py} &
Package export surface. &
Re-exports the main classes and functions. &
Makes imports cleaner in scripts and tests. \\
\bottomrule
\end{longtable}
\normalsize

\subsection{Important Script Files}
\scriptsize
\begin{longtable}{L{0.19\textwidth} L{0.27\textwidth} L{0.42\textwidth}}
\toprule
\textbf{File} & \textbf{Role} & \textbf{What it is used for} \\
\midrule
\endhead
\fname{build_qasper_processed.py} & Processed data builder & Generates canonical \fname{papers.jsonl}, \fname{questions.jsonl}, \fname{chunks.*.jsonl}, and \fname{manifest.json}. \\
\fname{evaluate_qasper_retrieval.py} & Retrieval runner & Runs retrieval-only experiments and writes retrieval JSON outputs. \\
\fname{evaluate_qasper_generation.py} & Generation runner & Runs retrieval + generation + evaluation end-to-end; now also supports citation verification. \\
\fname{qsub_qasper_retrieval.sh} & SCC batch wrapper & Runs retrieval jobs on SCC. \\
\fname{qsub_qasper_generation.sh} & SCC batch wrapper & Runs generation jobs on SCC with model, retrieval, and citation-verifier flags. \\
\fname{submit_proposal_ablation.sh} & Proposal batch chain & Submits the Phi A--E experiment chain. \\
\fname{submit_proposal_baselines.sh} & Baseline batch chain & Submits parametric-only, BM25-only, and random-retrieval baselines. \\
\fname{submit_full_validation_phi.sh} & Full-validation runner & Submits the strongest Phi settings on full validation. \\
\fname{submit_full_validation_ab.sh} & Full-validation runner & Submits full-validation A and B to complete the proposal matrix. \\
\fname{backfill_bertscore.py} & Metric backfill & Adds BERTScore to existing generation JSONs without rerunning generation. \\
\fname{backfill_nli_citation.py} & Metric backfill & Adds NLI citation metrics to existing generation JSONs. \\
\fname{analyze_generation_errors.py} & Error analysis & Groups generation failures into categories such as retrieval misses and unsupported citations. \\
\fname{compare_generation_runs.py} & Run comparison & Compares two saved generation runs on overlapping questions. \\
\fname{summarize_generation_runs.py} & Result packaging & Builds markdown tables from saved result JSONs. \\
\fname{prepare_human_eval_sample.py} & Human-eval preparation & Builds an enriched reviewer sample with gold answers, evidence, cited chunk text, and NLI claim hints. \\
\fname{create_human_eval_template.py} & Human-eval preparation & Turns a sample CSV into fillable reviewer templates. \\
\fname{aggregate_human_eval.py} & Human-eval aggregation & Merges completed reviewer CSVs into summary tables and JSON. \\
\fname{preview_qasper_chunking.py} & Debugging/sanity check & Lets the team inspect chunking behavior for a small paper subset. \\
\fname{prepare_shared_qasper.sh} & Dataset setup & Downloads and extracts public QASPER archives into the shared SCC directory. \\
\fname{pull_scc_results.sh} & Result archival & Safely copies SCC \fname{runs/} and \fname{logs/} locally without deleting the remote copies. \\
\fname{show_project_paths.py} & Debugging & Prints the resolved SCC/local path layout. \\
\fname{install_bertscore_overlay.sh} & SCC support utility & Installs a project-local overlay for BERTScore dependencies. \\
\fname{install_qwen3_overlay.sh} & SCC support utility & Installs a project-local transformers overlay needed for Qwen3. \\
\bottomrule
\end{longtable}
\normalsize

\subsection{Reports, Tests, and Artifacts}
\begin{itemize}
  \item \texttt{reports/proposal.tex}: original project proposal.
  \item \texttt{reports/midway\_report.tex}: ACL-formatted midway report.
  \item \texttt{reports/status\_report.tex}: internal running status report with the latest experiment state.
  \item \texttt{reports/references.bib}: shared bibliography.
  \item \texttt{reports/artifacts/final\_validation\_results.md}: final consolidated full-validation generation table.
  \item \texttt{reports/artifacts/citation\_tuning\_results.md}: tracked citation-tuning sample results.
  \item \texttt{reports/artifacts/human\_eval\_*}: reviewer sample, templates, instructions, and aggregated human-evaluation summary.
  \item \texttt{tests/test\_*.py}: unit and regression tests covering loading, chunking, retrieval, generation, NLI, configuration, and citation verification.
\end{itemize}

\section{Overall Workflow}
\subsection{From Input to Output}
\begin{enumerate}
  \item \textbf{Prepare raw QASPER data.} The project downloads and extracts the train/dev/test QASPER files.
  \item \textbf{Normalize papers and questions.} Raw papers, questions, answers, and evidence are parsed into a consistent internal schema.
  \item \textbf{Chunk papers.} Each paper is split into retrievable units using fixed-size, sliding-window, or section-aware chunking.
  \item \textbf{Write canonical processed artifacts.} These become the shared interface for all later experiments.
  \item \textbf{Run retrieval.} For each question, the system retrieves top chunks using BM25, dense retrieval, hybrid fusion, or reranking.
  \item \textbf{Build a prompt from retrieved chunks.} Retrieved chunks are compressed into the prompt context and labeled so the model can cite them.
  \item \textbf{Generate an answer.} The model produces either a baseline answer or a citation-forcing answer.
  \item \textbf{Normalize and parse the output.} The system extracts citations, cleans the answer, and normalizes cases like \texttt{yes}, \texttt{no}, and \texttt{unanswerable}.
  \item \textbf{Verify citations (new step).} In verifier mode, cited claims are checked with NLI. Unsupported cited outputs can be collapsed to \texttt{unanswerable}.
  \item \textbf{Evaluate.} Retrieval metrics, token F1, BERTScore, support rate, citation precision, and NLI citation precision are computed.
  \item \textbf{Archive and summarize.} SCC results are copied locally, summarized into markdown tables, and used for reports and human evaluation.
\end{enumerate}

\section{Results and Outputs}
\subsection{Retrieval Results}
\begin{table}[h]
\centering
\small
\begin{tabular}{lccc}
\toprule
\textbf{Method} & \textbf{Recall@5} & \textbf{MRR@10} & \textbf{Evidence Hit} \\
\midrule
BM25 & 0.5769 & 0.4110 & 0.7152 \\
Dense & 0.5886 & 0.4820 & 0.7238 \\
Hybrid & 0.6313 & 0.5097 & 0.7778 \\
Hybrid + rerank & 0.6410 & 0.5027 & 0.7702 \\
\bottomrule
\end{tabular}
\caption{Validation retrieval results on section-aware chunks.}
\end{table}

Interpretation:
\begin{itemize}
  \item Hybrid retrieval is the best balanced retriever.
  \item Reranking improves recall slightly, but that gain does not automatically translate into the best downstream generation quality.
\end{itemize}

\subsection{Full-Validation Generation Results}
\begin{table}[h]
\centering
\scriptsize
\begin{tabular}{lcccccccc}
\toprule
\textbf{Cond.} & \textbf{Chunking} & \textbf{Retrieval} & \textbf{Prompt} & \textbf{F1} & \textbf{BERT} & \textbf{Halluc.} & \textbf{Support} & \textbf{NLI Cite} \\
\midrule
A & fixed & dense & baseline & 0.3059 & 0.8791 & 0.5134 & 0.6886 & 0.0000 \\
B & fixed & hybrid & baseline & 0.3063 & 0.8782 & 0.5104 & 0.6716 & 0.0000 \\
C & section & hybrid & baseline & 0.3197 & 0.8815 & 0.4527 & 0.7741 & 0.0000 \\
D & section & hybrid\_rerank & baseline & \textbf{0.3238} & 0.8817 & \textbf{0.4478} & 0.7652 & 0.0000 \\
E & section & hybrid\_rerank & citation & 0.3215 & \textbf{0.8819} & 0.5463 & 0.5960 & 0.1416 \\
P & section & none & baseline & 0.1350 & 0.8329 & 0.8657 & 0.1343 & 0.0000 \\
M & section & BM25 & baseline & 0.3207 & 0.8813 & 0.4716 & 0.7254 & 0.0000 \\
\bottomrule
\end{tabular}
\caption{Main full-validation Phi results.}
\end{table}

What these mean:
\begin{itemize}
  \item \textbf{D is the best overall Phi setting.} Section-aware chunking plus hybrid retrieval with reranking and baseline prompting gives the best answer-quality and grounding balance.
  \item \textbf{E is good at producing citations, but not at trustworthy attribution.} It cites very often, but many cited claims are not actually supported.
  \item \textbf{P is much worse.} This shows retrieval is genuinely important; the model alone is not enough.
  \item \textbf{Section-aware chunking matters.} The jump from A/B to C/D shows that structural chunking helps more than simply changing retriever type on fixed chunks.
\end{itemize}

\subsection{Human Evaluation}
\begin{table}[h]
\centering
\small
\begin{tabular}{lccc}
\toprule
\textbf{System} & \textbf{Correctness} & \textbf{Grounding} & \textbf{Citation Quality} \\
\midrule
Baseline & 1.021 & 1.062 & -- \\
Citation & 0.875 & 0.708 & 0.583 \\
\bottomrule
\end{tabular}
\caption{Mean human-evaluation scores over 24 questions and 2 reviewers.}
\end{table}

Additional human-eval findings:
\begin{itemize}
  \item 24 questions were reviewed by 2 reviewers.
  \item Exact preference agreement was 83.3\%.
  \item Majority preference by question was split: 9 baseline, 9 citation, 6 tie.
\end{itemize}

Why this is significant:
\begin{itemize}
  \item Automatic metrics say baseline is stronger on correctness and grounding.
  \item Human judgments agree baseline is generally safer, but citation answers are not uniformly disliked.
  \item This means citation forcing is not useless, but it is not yet reliable enough to replace baseline as the main answer mode.
\end{itemize}

\section{Important Concepts and Design Decisions}
\begin{itemize}
  \item \textbf{Deterministic processed artifacts:} the team uses a shared \texttt{processed/} interface instead of rebuilding chunks differently on each machine.
  \item \textbf{Custom chunking instead of LangChain:} this keeps chunk IDs, metadata, and chunk boundaries stable and reproducible.
  \item \textbf{Hybrid retrieval:} combines sparse and dense retrieval using reciprocal-rank fusion.
  \item \textbf{Question-aware context compression:} retrieved chunks are compressed to relevant sentences before prompting, to fit SCC memory limits and reduce irrelevant context.
  \item \textbf{Two answer modes:} baseline prompting prioritizes answer quality, while citation forcing prioritizes explicit attribution.
  \item \textbf{Multi-metric evaluation:} token F1, BERTScore, support rate, hallucination proxy, citation precision, and NLI citation precision each capture different failure modes.
  \item \textbf{Move from prompt-only tuning to verification:} the new citation verifier is a deliberate shift from ``tell the model to cite'' toward ``check whether the cited claim is actually supported.''
\end{itemize}

\section{Limitations and Possible Improvements}
\subsection{Current Limitations}
\begin{itemize}
  \item Citation forcing still produces many unsupported cited claims.
  \item Token F1 is harsh and under-rewards semantically correct paraphrases.
  \item Human evaluation is still small.
  \item SCC scheduler/auth instability slowed iteration and some runs had to be retried.
  \item Some optional proposal stretch items remain secondary rather than central.
\end{itemize}

\subsection{Best Improvements}
\begin{itemize}
  \item Run the new verifier-gated citation mode on SCC and compare it directly to the current citation baseline.
  \item If the verifier helps, run it on full validation.
  \item Push citation mode toward a more explicit \emph{generate $\rightarrow$ verify $\rightarrow$ cite} pipeline.
  \item Add a second faithfulness metric only if it adds signal beyond the current NLI check.
  \item Keep stretch-model experiments secondary until the final writeup is locked.
\end{itemize}

\section{Simple Layman Explanation}
Imagine giving a long scientific paper to a student and asking, ``What baseline did they compare against?'' or ``Did they only test English?'' This project builds a system that tries to answer those questions automatically.

It works in three stages:
\begin{enumerate}
  \item it cuts the paper into smaller pieces,
  \item it searches those pieces to find the most relevant ones,
  \item it uses a language model to answer the question based only on those pieces.
\end{enumerate}

We tested different ways of cutting the paper up, different ways of searching for the right evidence, and different ways of forcing the model to show where its answer came from.

The main lesson is:
\begin{itemize}
  \item getting the right evidence matters a lot,
  \item the best current version answers best when allowed to answer normally,
  \item forcing the system to always add citations makes it cite more often, but those citations are not always truly trustworthy yet.
\end{itemize}

So in plain language, the project is a \textbf{research-paper question answering assistant} that already works reasonably well, but still needs a stronger evidence-checking step before its citations can be trusted.

\end{document}
